\section{A second sample of the diffusion model}
\label{sec:s-samples}
The diffusion row of the main text uses one sample per region.
A second sample on the same people, with every seed changed, lands at the same position (Table~\ref{tab:s-samples}): the paired difference is \mfSamplesCelebaDiffOne{} in CelebA and \mfSamplesImdbDiffOne{} in IMDb-WIKI, with intervals that include zero.
The pixelwise mean of the two samples is closer to the target than either of them, by \mfSamplesCelebaDiffMean{} and \mfSamplesImdbDiffMean{}, although it is no fill the model would produce.
This is the averaging reward of the main text, measured on the outputs of a generative model.
The second sample was added after the first was scored and is exploratory.

\begin{table}[htbp]
\centering
\caption{\textbf{Two samples of the diffusion model and their mean.}
\acs{MSE} on the random-person scale at 1/8 of the object, on the same people (column People) and windows.
Sample 0 is the reported run; differences are paired, with 95\% bootstrap intervals; the last two columns give the norm of the individual detail relative to the target's (median).
Averaging two samples moves the score toward the target while the fill keeps most of its detail.}
\label{tab:s-samples}
\small
\setlength{\tabcolsep}{4pt}
\begin{tabular}{@{}lrrrrllrr@{}}
\toprule
Set & People & Sample 0 & Sample 1 & Mean of 2 & Sample 1 $-$ sample 0 & Mean $-$ sample 0 & Detail, one & Detail, mean \\
\midrule
% Generated by paper/cvpr2027/make_model_fills.py from results/model_fills_v1/model_fills_sd_samples.json
CelebA & 58 & 0.130 & 0.132 & 0.102 & $+$0.002 [$-$0.004, 0.008] & $-$0.028 [$-$0.032, $-$0.024] & 0.95 & 0.92 \\
IMDb-WIKI & 43 & 0.136 & 0.138 & 0.107 & $+$0.002 [$-$0.004, 0.008] & $-$0.028 [$-$0.033, $-$0.024] & 0.95 & 0.93 \\
\bottomrule
\end{tabular}
\end{table}
% Source: results/model_fills_v1/model_fills_sd_samples.json (protocol amendment 3); rows by make_model_fills.py.
